\chapter[My second proper chapter]{My second proper chapter}  
\label{ch:first_model}
\section{An annoyingly long table}
Sometimes, you need a particularly annoyingly long table. Like this one which is in my thesis.
\begin{xltabular}{\linewidth}{lLl}
\caption[Summary of microphysical processes in CASIM]{A summary of the microphysical processes included in CASIM. Condensation and evaporation of cloud water is controlled outside CASIM by the bi-modal cloud fraction scheme of vWB21.}
\label{tab:CASIM_processes} \\

\hline \hline \multicolumn{1}{l}{\textbf{Process}} & \textbf{Description} & \multicolumn{1}{l}{\textbf{Hydrometeor}} \\ \hline \hline 
\endfirsthead

\multicolumn{3}{r}%
{\tablename\ \thetable{} -- continued from previous page} \\
\hline \hline \multicolumn{1}{l}{\textbf{Process}} & \textbf{Description} & \multicolumn{1}{l}{\textbf{Hydrometeor}}  \\ \hline \hline 
\endhead

\hline \multicolumn{3}{r}{{\emph{Continued on next page\dots}}} \\ \hline
\endfoot

\hline
\endlastfoot

Accretion &
  Collection of a (usually) smaller hydrometeor by a (usually) larger hydrometeor. & 
  \begin{tabular}[t]{@{}l@{}}cloud water by all others;\\ ice by all but cloud water;\\ rain by snow and graupel;\\ snow by graupel\end{tabular} \\ \hline
Activation &
  Activation of a cloud condensation nucleus into a droplet. By default, this is performed according to the ARG00 scheme. &
  cloud water only \\ \hline
Aggregation &
  Self-collection of precipitation-size hydrometeors of a single species. &
  rain and snow  \\ \hline
Autoconversion &
  Self-collection of a smaller species into the larger species of the same phase. &
  \begin{tabular}[t]{@{}l@{}}cloud water to rain;\\ ice to snow\end{tabular} \\ \hline
Deposition &
  Growth of ice particles by vapour deposition when air is supersaturated with respect to ice. Deposition ice nucleation is not included in this. &
  all frozen species \\ \hline
Evaporation &
  Evaporation rate of raindrops when entering air subsaturated with respect to water. &
  rain only  \\ \hline
\begin{tabular}[t]{@{}l@{}}Ice nucleation\\ (homogeneous)\end{tabular} &
  Freezing of liquid droplets at temperatures below \qty{-38}{\degreeCelsius}. &
  \begin{tabular}[t]{@{}l@{}}cloud water to ice;\\ rain to graupel\end{tabular} \\ \hline
\begin{tabular}[t]{@{}l@{}}Ice nucleation\\ (heterogeneous)\end{tabular} &
  Fully described elsewhere. By default, DM15 is used when aerosol is processed and C86 if not.
  
  Additionally, B53 is used for the heterogenous freezing of raindrops, though this is contained within the homogeneous nucleation subroutine. &
  cloud water to ice \\ \hline
Melting &
  Melting of frozen species to become rain. Instantaneous for ice, but not for snow and graupel, which adhere to a parameterised melting rate. &
  \begin{tabular}[t]{@{}l@{}}ice to rain;\\ snow to rain;\\ graupel to rain\end{tabular} \\ \hline
\begin{tabular}[t]{@{}l@{}}Secondary ice\\ production\end{tabular} &
  Increase in ice crystal concentration through the Hallett-Mossop process between \qty{-2.5}{\degreeCelsius} and \qty{-7.5}{\degreeCelsius}. The production rate is controlled by snow and graupel production rates. &
  all frozen species  \\ \hline
Sedimentation &
  Species falling into lower model levels under gravity. &
  all species \\ \hline
Sublimation &
  Sublimation of frozen hydrometeors when entering air subsaturated with respect to ice. &
  all frozen species \\ \hline \hline
\end{xltabular}